% ======================================================================
% Appendix D. Tauberian smoothing removal
% PATH: src/appendices/D-tauberian-smoothing-removal.tex
% ======================================================================

\appendix
\chapter{Tauberian smoothing removal and sharp counting recovery}
\label{app:D-tauberian}

\section{Purpose and scope}
\label{sec:D.purpose}

In several places of the main text, notably in Chapter~\ref{chap:09-local-weyl}, the primary spectral information is
extracted from \emph{smoothed} counting functionals rather than from sharp cutoffs.
This is natural from the perspective of trace formulae: admissible test functions and microlocal partitions yield
stable trace identities for mollified spectral measures, whereas sharp spectral projectors are discontinuous objects
whose analysis typically requires additional Tauberian input.

The purpose of this appendix is to record a self-contained Tauberian mechanism that upgrades smoothed trace
information to sharp counting statements, with explicit and stable remainder bookkeeping.
We isolate the analytic contribution of smoothing removal from the geometric contribution encoded in the trace
formula remainder ledger.

Throughout, we work at the level of general monotone measures and functions; the spectral applications follow by
specialization.

\section{Smoothing kernels and basic conventions}
\label{sec:D.kernels}

Let $w\in L^1(\R)$ satisfy the standard mollifier hypotheses:
\begin{equation}
\label{eq:D.mollifier.assumptions}
w\ge 0,
\qquad
\int_{\R} w(t)\,dt = 1,
\qquad
w(t)=w(-t).
\end{equation}
For $\delta>0$ define the rescaled kernel
\begin{equation}
\label{eq:D.wdelta}
w_\delta(t):=\delta^{-1}w(t/\delta).
\end{equation}

We write convolution as
\[
(f*w_\delta)(x)=\int_{\R} f(x-t)\,w_\delta(t)\,dt,
\]
whenever the integral is well-defined.

In trace formula applications $w$ is typically Schwartz or compactly supported smooth; however, for the basic
Tauberian inequalities only \eqref{eq:D.mollifier.assumptions} is needed.

\section{Monotone smoothing inequalities}
\label{sec:D.monotone}

\subsection{The elementary bracketing lemma}
\label{subsec:D.bracketing}

Let $F:\R\to\R$ be nondecreasing and locally bounded.
The fundamental observation is that convolution with a nonnegative probability kernel cannot increase the function
beyond its boundary layer.

\begin{lemma}[Monotone bracketing]
\label{lem:D.monotone-bracketing}
Let $F$ be nondecreasing and let $w_\delta$ satisfy \eqref{eq:D.mollifier.assumptions}.
Then for all $x\in\R$,
\begin{equation}
\label{eq:D.bracket}
F(x-\delta)
\le
(F*w_\delta)(x)
\le
F(x+\delta).
\end{equation}
\end{lemma}

\begin{proof}
Since $w_\delta$ is supported in $\R$ and integrates to $1$, one has
\[
(F*w_\delta)(x)
=
\int_{\R} F(x-t)\,w_\delta(t)\,dt.
\]
If $|t|\le \delta$ then $x-\delta\le x-t\le x+\delta$, and monotonicity implies
$F(x-\delta)\le F(x-t)\le F(x+\delta)$.
Integrating against the probability measure $w_\delta(t)\,dt$ yields \eqref{eq:D.bracket}.
\end{proof}

\subsection{A quantitative smoothing deviation bound}
\label{subsec:D.deviation}

The bracketing immediately implies a stable estimate for the difference between $F$ and its smoothing.

\begin{corollary}[Deviation controlled by boundary increment]
\label{cor:D.boundary-increment}
Under the hypotheses of Lemma~\ref{lem:D.monotone-bracketing},
\begin{equation}
\label{eq:D.deviation}
\bigl|(F*w_\delta)(x)-F(x)\bigr|
\le
F(x+\delta)-F(x-\delta).
\end{equation}
\end{corollary}

\begin{proof}
From \eqref{eq:D.bracket} we have
\[
(F*w_\delta)(x)-F(x)
\le
F(x+\delta)-F(x),
\qquad
F(x)-(F*w_\delta)(x)
\le
F(x)-F(x-\delta).
\]
Summing the two inequalities yields \eqref{eq:D.deviation}.
\end{proof}

\section{Spectral measures and smoothed counting functions}
\label{sec:D.spectral}

Let $\mu$ be a positive Radon measure on $[0,\infty)$.
Define the sharp counting function
\begin{equation}
\label{eq:D.sharp-count}
N(r):=\mu([0,r]),
\qquad r\ge 0.
\end{equation}
Assume $N$ is finite for each $r$.

Define the smoothed counting function
\begin{equation}
\label{eq:D.smoothed-count}
N_\delta(r)
:=
(N*w_\delta)(r)
=
\int_{\R} N(r-t)\,w_\delta(t)\,dt,
\qquad r\ge 0,
\end{equation}
where $N(r)=0$ for $r<0$.

In the spectral setting of Chapter~\ref{chap:09-local-weyl}, one has
\[
\mu=\sum_j \delta_{r_j},
\qquad
N(r)=\#\{j:\ r_j\le r\},
\qquad
N_\delta(r)=\sum_j w_\delta(r-r_j).
\]

The Tauberian mechanism consists of proving asymptotics for $N_\delta(r)$ and then controlling the difference
$N_\delta(r)-N(r)$ by boundary increments.

\section{Tauberian upgrade principle: a stable formulation}
\label{sec:D.tauberian-upgrade}

\subsection{A general upgrade lemma}
\label{subsec:D.upgrade-lemma}

The following statement is the central abstract tool.

\begin{lemma}[Tauberian upgrade: smoothed to sharp]
\label{lem:D.upgrade}
Let $N$ be a nondecreasing function on $[0,\infty)$.
Assume that for some $\alpha>0$ and some constant $C>0$ one has a smoothed asymptotic
\begin{equation}
\label{eq:D.smoothed-asymptotic}
N_\delta(r)
=
C r^\alpha
+
R_\delta(r),
\qquad r\ge 1,
\end{equation}
where $R_\delta(r)$ is an error term.

Then for all $r\ge 1$,
\begin{equation}
\label{eq:D.sharp-recovery}
N(r)
=
C r^\alpha
+
R_\delta(r)
+
O\!\bigl(N(r+\delta)-N(r-\delta)\bigr).
\end{equation}
\end{lemma}

\begin{proof}
By Corollary~\ref{cor:D.boundary-increment},
\[
|N_\delta(r)-N(r)|
\le
N(r+\delta)-N(r-\delta).
\]
Inserting \eqref{eq:D.smoothed-asymptotic} yields \eqref{eq:D.sharp-recovery}.
\end{proof}

Lemma~\ref{lem:D.upgrade} is completely robust: it requires no smoothness, only monotonicity.
All problem-specific information is concentrated in the boundary increment term.

\subsection{A scale selection criterion}
\label{subsec:D.scale-selection}

To extract an asymptotic from \eqref{eq:D.sharp-recovery}, one must select $\delta=\delta(r)$ such that
the increment term is of lower order than $r^\alpha$.

A canonical sufficient condition is:
\begin{equation}
\label{eq:D.increment-negligible}
N(r+\delta(r))-N(r-\delta(r))=o(r^\alpha),
\qquad r\to\infty.
\end{equation}

In spectral Weyl regimes, such control follows from a crude upper bound on $N$ itself.

\section{Increment control under polynomial growth}
\label{sec:D.polynomial-growth}

\subsection{A deterministic increment bound}
\label{subsec:D.increment-bound}

Assume a global polynomial bound on $N$:
\begin{equation}
\label{eq:D.polynomial-bound}
N(r)\le C_0(1+r)^\alpha,
\qquad r\ge 0.
\end{equation}

Then for small $\delta$ one obtains a deterministic increment estimate.

\begin{lemma}[Polynomial growth increment bound]
\label{lem:D.poly-increment}
Assume \eqref{eq:D.polynomial-bound} holds for some $\alpha>0$.
Then for all $r\ge 1$ and all $\delta\in(0,1]$,
\begin{equation}
\label{eq:D.increment-poly}
N(r+\delta)-N(r-\delta)
\le
C_1\,\delta\,r^{\alpha-1},
\end{equation}
where $C_1$ depends only on $C_0$ and $\alpha$.
\end{lemma}

\begin{proof}
By monotonicity and \eqref{eq:D.polynomial-bound},
\[
N(r+\delta)-N(r-\delta)
\le
C_0\bigl((1+r+\delta)^\alpha-(1+r-\delta)^\alpha\bigr).
\]
The function $x\mapsto (1+x)^\alpha$ is $C^1$ on $[0,\infty)$ and satisfies
\[
(1+r+\delta)^\alpha-(1+r-\delta)^\alpha
\le
2\delta \sup_{|t|\le \delta} \alpha(1+r+t)^{\alpha-1}
\ll
\delta\,r^{\alpha-1}.
\]
This yields \eqref{eq:D.increment-poly}.
\end{proof}

\subsection{A canonical smoothing removal regime}
\label{subsec:D.canonical-regime}

If the smoothed law \eqref{eq:D.smoothed-asymptotic} holds uniformly for $\delta\in(0,1]$ and if
\eqref{eq:D.polynomial-bound} holds, then choosing $\delta(r)=r^{-\theta}$ with $\theta\in(0,1)$ yields
\[
N(r+\delta(r))-N(r-\delta(r))
\ll
r^{\alpha-1-\theta}
=
o(r^\alpha).
\]
Thus smoothing removal is valid with a power-saving margin.

\section{A two-component remainder decomposition}
\label{sec:D.two-component}

\subsection{Geometric remainder versus Tauberian remainder}
\label{subsec:D.geometric-vs-tauberian}

In the main text the smoothed asymptotic \eqref{eq:D.smoothed-asymptotic} arises from a trace formula identity.
It is conceptually essential to separate the error sources:

\begin{enumerate}
\item the \emph{trace remainder} (geometric, normalization, continuous-spectrum terms), controlled by the ledger;
\item the \emph{Tauberian remainder} (smoothing removal), controlled by the increment term.
\end{enumerate}

We formalize this separation as follows.

\begin{definition}[Tauberian remainder]
\label{def:D.tauberian-remainder}
Let $N$ be a counting function and $N_\delta=N*w_\delta$ its smoothing.
Define the Tauberian remainder at scale $\delta$ by
\begin{equation}
\label{eq:D.tauberian-remainder}
\mathcal{R}^{\mathrm{Tab}}(r;\delta)
:=
N_\delta(r)-N(r).
\end{equation}
\end{definition}

By Corollary~\ref{cor:D.boundary-increment} one has the deterministic bound
\begin{equation}
\label{eq:D.tab-bound}
\bigl|\mathcal{R}^{\mathrm{Tab}}(r;\delta)\bigr|
\le
N(r+\delta)-N(r-\delta).
\end{equation}

In particular, the Tauberian remainder is a purely analytic boundary-layer quantity and does not depend on the
geometric structure of the trace formula.

\subsection{A stable sharp Weyl recovery statement}
\label{subsec:D.sharp-weyl-recovery}

We now record a standard sharp recovery theorem in the form used in Chapter~\ref{chap:09-local-weyl}.

\begin{theorem}[Sharp recovery from smoothed Weyl asymptotics]
\label{thm:D.sharp-from-smoothed}
Let $N$ be a nondecreasing counting function.
Assume that for each $\delta\in(0,1]$ one has
\begin{equation}
\label{eq:D.assume-smoothed}
N_\delta(r)
=
C r^\alpha
+
R^{\mathrm{TF}}(r;\delta),
\qquad r\ge 1,
\end{equation}
with a trace-formula remainder bound
\begin{equation}
\label{eq:D.assume-TF-bound}
|R^{\mathrm{TF}}(r;\delta)|
\le
\mathcal{E}^{\mathrm{TF}}(r;\delta).
\end{equation}
Then for all $r\ge 1$,
\begin{equation}
\label{eq:D.sharp-final}
N(r)
=
C r^\alpha
+
O\!\bigl(\mathcal{E}^{\mathrm{TF}}(r;\delta)\bigr)
+
O\!\bigl(N(r+\delta)-N(r-\delta)\bigr).
\end{equation}
\end{theorem}

\begin{proof}
Combine Lemma~\ref{lem:D.upgrade} with the assumed bound \eqref{eq:D.assume-TF-bound}.
\end{proof}

\section{Weighted Tauberian smoothing removal}
\label{sec:D.weighted}

\subsection{Weighted counting functions}
\label{subsec:D.weighted-count}

In the microlocal Weyl laws one encounters weighted counting measures of the form
\[
\mu^A := \sum_j \langle A\varphi_j,\varphi_j\rangle\,\delta_{r_j},
\]
where $A\in\Psi^0$.
The corresponding counting function is
\[
N_A(r):=\mu^A([0,r])
=
\sum_{r_j\le r} \langle A\varphi_j,\varphi_j\rangle.
\]

If $A$ is bounded on $L^2$, then $N_A$ is of bounded variation but is not necessarily monotone.
Nevertheless, Tauberian removal remains possible under a uniform boundedness hypothesis.

\subsection{A bounded-variation upgrade lemma}
\label{subsec:D.bv}

We record the simplest form needed for the main text.

\begin{lemma}[Smoothing removal under bounded variation]
\label{lem:D.bv-upgrade}
Let $F$ be a function of bounded variation on $\R$.
Assume $w_\delta$ satisfies \eqref{eq:D.mollifier.assumptions}.
Then for all $x$,
\[
|(F*w_\delta)(x)-F(x)|
\le
\Var(F;[x-\delta,x+\delta]),
\]
where $\Var(F;I)$ denotes the total variation of $F$ on $I$.
\end{lemma}

\begin{proof}
Write $F$ as the difference of two monotone functions (Jordan decomposition) and apply
Corollary~\ref{cor:D.boundary-increment} to each component.
\end{proof}

In applications, $F=N_A$ and $\Var(F;I)$ is bounded by $\|A\|_{L^2\to L^2}$ times the number of eigenvalues in the
interval $I$.

\section{A minimal Tauberian template for the main text}
\label{sec:D.template}

For convenience, we summarize the Tauberian upgrade mechanism in a compact reusable template.

\begin{proposition}[Tauberian template]
\label{prop:D.template}
Let $N$ be a sharp counting function and let $N_\delta=N*w_\delta$ be its smoothing.
Assume the smoothed asymptotic
\[
N_\delta(r)=M(r)+R(r;\delta),
\]
where $M(r)$ is an explicit main term.
Then
\[
N(r)=M(r)+R(r;\delta)+O\!\bigl(N(r+\delta)-N(r-\delta)\bigr).
\]
If in addition $N(r)\ll r^\alpha$ and $\delta=r^{-\theta}$ with $\theta\in(0,1)$, then
\[
N(r+\delta)-N(r-\delta)=o(r^\alpha),
\]
and hence the sharp asymptotic $N(r)\sim M(r)$ follows provided $R(r;\delta)=o(r^\alpha)$.
\end{proposition}

\begin{proof}
Immediate from Lemma~\ref{lem:D.upgrade} and Lemma~\ref{lem:D.poly-increment}.
\end{proof}

\section*{Bibliographic notes}

The bracketing inequality of Lemma~\ref{lem:D.monotone-bracketing} is classical and appears in essentially all
Tauberian arguments for Weyl laws.
The upgrade mechanism stated in Theorem~\ref{thm:D.sharp-from-smoothed} is the standard route from smoothed trace
formula asymptotics to sharp counting statements, and may be viewed as a minimal analytic interface between
microlocal trace identities and sharp spectral distribution results.

% ======================================================================
% End of Appendix D
% ======================================================================
